\section{Kết luận}

Trong phần Kết luận, cần khái quát lại mục tiêu, phương pháp và những phát hiện quan trọng nhất của nghiên cứu mà không lặp lại toàn bộ kết quả đã trình bày. Trước hết, tóm lược cách tiếp cận phân tích dữ liệu và trắc lượng thư mục được sử dụng để khảo sát các ấn phẩm Khoa học Thông tin và Thư viện bị thu hồi trong giai đoạn 2021 - 2026. Tiếp theo, nhấn mạnh những phát hiện có ý nghĩa nhất, chẳng hạn xu hướng thu hồi theo thời gian, đặc điểm của các ấn phẩm và tạp chí liên quan, phân bố theo lĩnh vực hoặc chủ đề, đặc điểm hợp tác và các nguyên nhân thu hồi nổi bật. Sau đó, khái quát ý nghĩa của các phát hiện đối với việc nhận diện và hiểu rõ các mô hình thu hồi ấn phẩm, đồng thời nhấn mạnh giá trị của việc kết hợp phân tích dữ liệu với trắc lượng thư mục trong nghiên cứu liêm chính khoa học. Cuối cùng, kết luận cần nêu rõ phạm vi của các phát hiện và gợi mở hướng nghiên cứu tiếp theo, đặc biệt là mở rộng nguồn dữ liệu, theo dõi các ấn phẩm theo thời gian và phát triển các phương pháp chuẩn hóa để phân biệt xu hướng thu hồi thực sự với ảnh hưởng của các sự kiện thu hồi hàng loạt.